\subsection{Restriction, descent, and persistent obstruction germs}

Write $A=N/\gamma_3N$ and
$\operatorname{Sec}_L=H^1_{\rm cts}(G_L,A)$ for a number field
$L\subset\overline{\mathbb Q}$.  The action is the restriction of
the action based at the fixed tangent at $O$.
Our maps below use that actual action, and all equivalences of
sections are by conjugation in $A$.
This is the germ construction requested in CK C9, cote 149,
\S3, manuscript p.~33, scan 66, applied to the stated class-two
quotient; the Fermat comparison is \S51.16, folios
733v--736v, scans 172--178.

\begin{statement}{C.1}{Proved}{No arithmetic centralizer in the class-two geometric group}
For every number field $L$ and every continuous section
$s:G_L\to A\rtimes G_L$,
\[
                 C_A(s(G_L))=1.                         \tag{C.1}
\]
The assertion uses the whole $G_L$, in contrast to the subgroup
$\ker\chi$ in B.4; the comparison map between the two assertions
is the inclusion of subgroups, which reverses inclusions of their
centralizers.
\end{statement}
\begin{proof}
Let $h$ be in the displayed centralizer.  Its image in
$V=A^{\rm ab}$ is fixed by $G_L$, since inner conjugation is
trivial on abelianization.  Choose a finite prime $v$ of $L$
above two, with residue size $q=2^f$.  The eigenvalues of
arithmetic Frobenius on $W\otimes\mathbb Q_3$ have complex
absolute value $q$: they are $q$ times roots of unity from the
finite cusp permutation action.  On $T_3E\otimes\mathbb Q_3$
the two eigenvalues have absolute value $q^{1/2}$.
Here the necessary elliptic calculation is explicit:
$E(\mathbb F_2)$ has three points by
\texttt{bc\_checks.py}; the Frobenius $F_2$ has determinant
two by the Weil pairing, and
$\det(1-F_2)=\deg(1-F_2)=3$ because $1-F_2$ is separable
with precisely those three kernel points.
Thus its characteristic polynomial is $X^2+2$.
On the base change at $v$ the Frobenius is $F_2^f$,
which proves the asserted absolute values directly.
The characteristic polynomial of the extension $V$ is the
product of those on $W$ and $T$, so one is not an eigenvalue.
Consequently $V^{G_L}=0$, and $h$ lies in the central group
$\gamma_2N/\gamma_3N=\bigwedge^2V$.
Inner conjugation acts trivially there.  Its Frobenius
eigenvalues, as seen from the filtration induced by
$0\to W\to V\to T\to0$, have absolute values
$q^2,q^{3/2},q$.  Again none is one.  The lattice is
torsion-free, so its invariant subgroup is zero.  Hence $h=1$.
\end{proof}

\begin{statement}{C.2}{Proved}{All restriction maps are injective and finite Galois descent is exact}
For number fields $L\subset L'$, restriction is the injective
map of pointed sets
\[
 \operatorname{res}_{L'/L}:\operatorname{Sec}_L
                   \longrightarrow\operatorname{Sec}_{L'},
 \qquad [z]\longmapsto[z|_{G_{L'}}].                        \tag{C.2}
\]
For $L'/L$ finite Galois it identifies its source with precisely
the $\operatorname{Gal}(L'/L)$-fixed classes in its target.
Writing $\rho$ for the based action on $A$, the action by
$g\in G_L$ is represented by
\[
       (g\cdot z)(h)=\rho(g)z(g^{-1}hg),\quad h\in G_{L'}.
                                                               \tag{C.3}
\]
Thus the kernel of (C.2), as a pointed-set map, is its base
point, and its image in the invariant set is the entire set.
\end{statement}
\begin{proof}
Restriction preserves the continuous cocycle equation and
geometric conjugation, proving that (C.2) is defined.
For (C.3), conjugate the section graph by $(1,g)$ in
$A\rtimes G_L$.  This gives exactly the formula and proves
continuity and the cocycle equation.  When $g\in G_{L'}$,
the cocycle identity gives
\[
 \rho(g)z(g^{-1}hg)=z(g)^{-1}z(h)\rho(h)z(g),
\]
so the action is conjugation by $z(g)^{-1}$ and is trivial
on section classes.  Therefore it factors through the
specified finite Galois group.

For injectivity, conjugate two sections so that their
restrictions to $G_{L'}$ agree.  Choose a finite extension
$M/L$, Galois over $L$, with $M\supset L'$.
The restrictions agree on the normal subgroup $G_M$.
For each $g\in G_L$, the element $t(g)s(g)^{-1}\in A$
centralizes $s(G_M)$: both sections induce, by conjugation,
the same map $s(h)\mapsto s(ghg^{-1})$ there.
Statement C.1 gives $t(g)s(g)^{-1}=1$, proving equality
over $G_L$.

For the image assertion, let $s$ represent a fixed class
over $L'$.  Consider its normalizer
\[
       B=N_{A\rtimes G_L}(s(G_{L'})).
\]
It is a closed subgroup of a compact profinite group.
Invariance of the class under (C.3) says precisely that
each $g\in G_L$ is the projection of an element of $B$:
use the conjugator in $A$ comparing $g\cdot s$ with $s$.
The kernel of $B\to G_L$ equals
$C_A(s(G_{L'}))$.  Indeed, an element projecting to one
normalizes the graph only by preserving each of its
points, since there is exactly one graph point over
each $h\in G_{L'}$.
This kernel is one by C.1.  Hence $B\to G_L$ is a
continuous bijective homomorphism between compact
Hausdorff groups and has continuous inverse.  That
inverse is a section extending $s$, because its unique
element over $h\in G_{L'}$ is $s(h)$.
Injectivity already proved uniqueness up to geometric
conjugacy.
\end{proof}

\begin{statement}{C.3}{Proved}{The action and fields of definition of germs}
The set
\[
          \mathscr G=\varinjlim_{L\subset\overline{\mathbb Q}}
                                      \operatorname{Sec}_L
                                                               \tag{C.4}
\]
has a well-defined $G_{\mathbb Q}$-action obtained by
(C.3) together with conjugation of the field of definition.
Every transition map is injective.  Each germ has an open
stabilizer and a finite orbit.  For every number field $L$,
\[
       \operatorname{Sec}_L\ \xrightarrow{\sim}\
                                    \mathscr G^{G_L}.          \tag{C.5}
\]
In particular the image of $\operatorname{Sec}_{\mathbb Q}$
is exactly the fixed-germ set.  The fixed field of the
stabilizer of a germ is its smallest field of definition.
\end{statement}
\begin{proof}
The formula transports a cocycle on $G_L$ to a cocycle on
$G_{gL}$, and respects restriction and conjugation.
It therefore defines an action on the colimit.
Restriction injectivity is C.2.  A representative over
$L$ is fixed by $G_L$, since the same cocycle calculation
used after (C.3) makes inner field conjugation trivial
on its class.  The stabilizer is consequently open.

Take a germ fixed by $G_L$, choose a field where it is
represented, and enlarge it to a finite Galois extension
$M/L$.  For $g\in G_L$, the class and its translate are
both defined over $M$.  Their equality as germs means
equality after restriction to a further finite field;
injectivity in C.2 makes them equal already over $M$.
The invariant class descends by C.2 to $L$.  Uniqueness
is again injectivity.  This proves (C.5).

An open stabilizer $H\subset G_{\mathbb Q}$ is $G_{L_0}$
for a number field $L_0$.  Applying the same argument
over $L_0$ proves definability there.  Every other field
of definition has its Galois group inside $H$, so it
contains $L_0$.  This proves the final assertion.
\end{proof}

\begin{statement}{C.4}{Refuted}{The R2 obstruction does not disappear as a germ}
Let $c=9\kappa(3)\cup\beta$ be the nonzero R2 class
specified by LM7 (E.32f), so that in
$H^2_{\rm cts}(G_{\mathbb Q},\mathbb Z_3(2))=
\mathcal R/d\mathcal R$ it is represented by
\[
 z_{q_j}=q_j-1\quad(j\ge3),\qquad z_q=0
                       \text{ at the other primes}.          \tag{C.6}
\]
For every finite number field $L$,
$\operatorname{res}_{L/\mathbb Q}(c)\ne0$.
In fact these restrictions all have infinite additive order.
Consequently $c$ gives a nonzero element of
$\varinjlim_L H^2_{\rm cts}(G_L,\mathbb Z_3(2))$.
The vanishing assertion in the first sentence of Task C.4
is false.
\end{statement}
\begin{proof}
For a nonzero integer $m$, the unique coordinatewise
solution of $dy=mz$ has value $m$ at every $q_j$ and zero
elsewhere, because multiplication by each nonzero
$q-1$ is injective in $\mathbb Z_3$.
This family does not tend to zero $3$-adically.  Thus
$mz\notin d\mathcal R$ and $mc\ne0$ for every nonzero
integer $m$.

We spell out the transfer used for this continuous
cohomology.  Set $G=G_{\mathbb Q}$, $H=G_L$ and
$M=\mathbb Z_3(2)$.  Use homogeneous continuous
equivariant cochains, with alternating deletion
differential.  Choose representatives $T$ for the
finitely many left cosets $H\backslash G$, and write
$r(g)=gt(g)^{-1}\in H$ for $g\in Ht(g)$.
For an $H$-cochain $c$ define the $G$-cochain
\[
 (\mathcal Tc)(g_0,\ldots,g_n)=
 \sum_{t\in T}t^{-1}c(r(tg_0),\ldots,r(tg_n)).              \tag{C.7a}
\]
The coset permutation $tg=h(t,g)t'$ gives
$t^{-1}h(t,g)=g(t')^{-1}$, proving equivariance;
the formula commutes with the differential.
All maps are continuous because the transversal is
finite and each coset is open.  Put
$a_t(g)=t^{-1}r(tg)$.  For a degree-$n$ $G$-cochain $f$,
the prism operator is
\[
 (Pf)(g_0,\ldots,g_{n-1})=
 \sum_{t\in T}\sum_{i=0}^{n-1}(-1)^i
 f(a_tg_0,\ldots,a_tg_i,g_i,\ldots,g_{n-1}).              \tag{C.7b}
\]
The same coset permutation makes the sum equivariant.
In $dP+Pd$ the internal mixed faces cancel in pairs;
the two end faces are $[G:H]f$ and
$-\mathcal T(\operatorname{res}f)$.
Thus the formulas prove directly on continuous
$\mathbb Z_3(2)$ cochains that
\[
 \operatorname{cor}_{L/\mathbb Q}
 \operatorname{res}_{L/\mathbb Q}=[L:\mathbb Q].              \tag{C.7}
\]
Equivalently one can apply the same compatible formulas
to every finite coefficient module; their cochain
homotopies commute with reduction, so the assertion
holds on the derived inverse limit as well as on
finite-coefficient cohomology.  This includes its
$\varprojlim^1$ subgroup.

A zero restriction would make $[L:\mathbb Q]c=0$,
contrary to the first paragraph.  A torsion restriction
would, by the same calculation, make a nonzero integer
multiple of $c$ zero.  Thus it has infinite order.
An element maps to zero in a filtered colimit of groups
exactly when it maps to zero at some later stage;
no such stage exists here.
\end{proof}

\begin{statement}{C.5}{Proved}{Restriction is injective on the Tate-twist phantom subgroup}
For $\mathbb Z_3(2)$ the kernel of the finite-coefficient
comparison map in R2 is torsion-free, and restriction
from $\mathbb Q$ to any finite number field is injective
on this kernel.  This extends C.4 from its particular
cup-product witness to every nonzero class in that
kernel.
\end{statement}
\begin{proof}
The kernel is
\[
 \frac{\{z\in\mathcal R:z_q\in(q-1)\mathbb Z_3\text{ for all }q\}}
      {d\mathcal R}.
\]
Write $y_q=z_q/(q-1)\in\mathbb Z_3$.  A nonzero integer
multiple $mz$ lies in $d\mathcal R$ exactly when
$(my_q)_q$ is restricted.  Multiplication by a fixed
nonzero integer changes valuations by the fixed
integer $v_3(m)$, so this condition is equivalent to
$(y_q)_q$ being restricted.  This proves torsion-freeness.
Equation (C.7) then proves injectivity of restriction
on this subgroup.  No closure is taken in the
denominator in this argument.
\end{proof}

\begin{statement}{C.7}{Proved}{The intrinsic phantom obstruction also survives every finite extension}
The nonzero obstruction in $H^2_{\rm cts}(G_{\mathbb Q},\mathcal K)$
constructed in A.4 remains nonzero over every finite number
field and in the filtered colimit over all number fields.
More generally, restriction to every finite number field is
injective on the phantom subgroup for $\mathcal K$ computed
in A.2.
\end{statement}
\begin{proof}
In the residue presentation of A.2, every
$D_q=\operatorname{Frob}_q-1$ on $\mathcal K(-1)$ is
injective.  A locally zero class has a unique integral
correction $y_q=D_q^{-1}z_q$ at each prime.
A nonzero integer multiple of its class vanishes
precisely when $(my_q)_q$ is restricted, which is
equivalent to $(y_q)_q$ being restricted.
Thus this phantom subgroup is torsion-free.
The cochain transfer and prism formulas (C.7a)--(C.7b)
apply verbatim with $M=\mathcal K$; they prove
corestriction after restriction equals multiplication
by the finite extension degree.  Such multiplication
is injective on the phantom subgroup.
This proves the assertion and completes both
persistence decisions asked for in Task C.4.
\end{proof}

\begin{statement}{C.10}{Proved}{The nine disjoint cusp germ packets and their Galois action}
Define the abelian germ parameter group
\[
 \mathscr K=\varinjlim_{L\subset\overline{\mathbb Q}}
       H^1_{\rm cts}(G_L,\mathbb Z_3(1))
   =\varinjlim_L\varprojlim_n L^\times/L^{\times3^n}.       \tag{C.10.1}
\]
All transition maps are restriction.  Choose at each cusp
$d\in D(\overline{\mathbb Q})$ the unit tangent for the parameter
$x/z$ when $x=0$, $y/z$ when $y=0$, and $z/x$ when $z=0$.
These choices are $G_{\mathbb Q}$-equivariant.
The map
\[
 \coprod_{d\in D(\overline{\mathbb Q})}\mathscr K
       \longrightarrow\mathscr G,
 \qquad (d,u)\longmapsto[s_{d,1}\text{ twisted by }u]
                                                               \tag{C.10.2}
\]
is injective, and its image is exactly the cuspidal germ set.
There are nine packets.  The $G_{\mathbb Q}$-action in these
coordinates is
\[
                         g(d,u)=(gd,gu),                     \tag{C.10.3}
\]
where $gu$ transports a Kummer class over $L$ to its class over
$gL$ by the field isomorphism $g:L\to gL$.
The parameter group has injective restriction maps, is divisible,
and has torsion subgroup canonically $\mu_{3^\infty}$, with its
usual Galois action.  In particular, distinct torsion parameters
and distinct cusp packets do not merge in the germ colimit.
\end{statement}
\begin{proof}
The unit tangent gives a section of the decomposition extension
at $d$.  Over a number field $L$ containing $d$, this extension
is the semidirect product $\mathbb Z_3(1)\rtimes G_L$.
Its continuous sections modulo inertia conjugacy are precisely
$H^1(G_L,\mathbb Z_3(1))$.  The decomposition homomorphism into
$N\rtimes G_L$ and the continuous quotient $N\to A$ define
the map in (C.10.2).  Passing to a larger finite field restricts
the same section, so the map is defined on (C.10.1).
The fields containing $d$ form a cofinal subsystem of the fields
in that definition.  Every cuspidal germ is, by definition, a
section of one such decomposition extension, proving the image
assertion once injectivity is checked.

First note that $H^0(G_L,\mathbb Z_3(1))=0$ for every finite
number field $L$.  If a one-cocycle becomes a coboundary over
a finite extension, subtract that coboundary and enlarge to a
finite Galois extension $M/L$ so that the cocycle vanishes on
$G_M$.  For $g\in G_L$, its value is fixed by $G_M$, by applying
the cocycle identity to $gG_Mg^{-1}=G_M$.  The displayed
invariant group is zero.  The cocycle therefore vanishes on
$G_L$, which proves injectivity of every restriction in
(C.10.1).

Enlarge a common field of definition to contain all nine cusps.
The boundary module is
$W=\mathbb Z_3(1)[D]/\mathbb Z_3(1)\sum_{e\in D}e_e$.
For a given $d$, the map $e_d:\mathbb Z_3(1)\to W$ is a
split $G_L$-equivariant inclusion: omit any other cusp from
the boundary basis.  The map $H^1(G_L,W)\to H^1(G_L,V)$
is injective because $H^0(G_L,T_3E)=0$, as proved in C.1.
Thus $e_d:H^1(G_L,\mathbb Z_3(1))\to H^1(G_L,V)$ is
injective.  The abelian image of the packet is exactly
\[
                   \alpha_d+e_d H^1(G_L,\mathbb Z_3(1)),     \tag{C.10.4}
\]
where $\alpha_d$ is the Kummer class of the framed tangential
divisor $d-O$.  This follows from the same division-torsor
comparison used in B.1, now over $L$.
Consequently two parameters in one packet cannot become
equal as abelian germs, and hence cannot become equal in
$\mathscr G$.

Filling the boundary gives the continuous equivariant quotient
$V\twoheadrightarrow T_3E$.  It sends (C.10.4) to
$\kappa_L(d-O)$.  For $d\ne e$, the difference $d-e$ is a
nonzero point of $E[3]$: the flex tangent relation is
$3d\sim3O$, and the nine distinct flexes exhaust $E[3]$.
The integral Kummer class $\kappa_L(d-e)$ remains nonzero
over every finite extension of $L$.  Indeed the
Mordell--Weil group of that extension is finitely generated.
A nonzero three-primary torsion element in such a group is
not divisible by $3^n$ for all $n$, so has nonzero image in
its completion and in $H^1(T_3E)$ by the Kummer exact sequence.
This argument uses the entire integral Kummer class; its
modulo-three image alone can vanish after adjoining ninth
division points.  It proves that the nine packets remain
pairwise disjoint in the colimit.

Our parameters and tangent value one are defined by the
coordinate functions over $\mathbb Q$.  Field conjugation
therefore takes the chosen tangent at $d$ to the chosen
tangent at $gd$.  Transporting the decomposition inclusion
conjugates its geometric path, and geometric conjugacy is
exactly the equivalence used in $\mathscr G$.
Kummer functoriality proves (C.10.3), with no additional
translation term.

For completeness, a finite-level Kummer class modulo three
is killed by some finite extension: use the finite image of
its affine cocycle representation.  The long exact sequence
of $0\to\mathbb Z_3(1)\xrightarrow{3}\mathbb Z_3(1)
\to\mu_3\to0$ then shows that the corresponding germ is
divisible by three.  Multiplication by an integer prime to
three is already invertible on each completed Kummer group.
Thus $\mathscr K$ is divisible.
At a fixed finite field, its completed Kummer group's torsion
is precisely $\mu_{3^\infty}(L)$.  To see this directly, use
the divisor sequence from $L^\times$ to its principal
fractional ideals.  The latter group is free abelian and the
sequence splits as abelian groups; the unit group is a finite
roots-of-unity group plus a free abelian group by Dirichlet's
theorem.  The three-adic completion of a free abelian group
is a subgroup of a product of copies of $\mathbb Z_3$, hence
torsion-free.  Only the three-primary roots of unity remain
as torsion.  Since all restrictions in (C.10.1) are injective,
passing to the colimit introduces no further torsion and
retains all roots of unity of three-power order.  This proves
the final assertions.
\end{proof}

\begin{statement}{C.11}{Refuted}{The pure-residue computation over the rationals does not extend verbatim to every number field}
Let
\[
       L=\mathbb Q(\zeta_3,\sqrt{10}),\qquad
       S_L=\{v:v\mid3\}\cup\{v:v\mid\infty\}.
\]
For the actual rank-forty-four module $\mathcal K$,
\[
                   H^2_{\rm cts}(G_{L,S_L},\mathcal K)\ne0.
                                                               \tag{C.11.1}
\]
Its inflation into $H^2_{\rm cts}(G_L,\mathcal K)$ is
injective, and every one of its classes has zero
residue outside $S_L$.
Consequently the pure-residue isomorphism (A.1.3)
cannot be transported to this $L$ by replacing the
rational primes with its finite primes.  The
finite-$S_L$ complex is a necessary part of Task C.1.
\end{statement}
\begin{proof}
The polynomial $X^2-10$ has two simple roots modulo
three and hence two roots in $\mathbb Q_3$ by Hensel's
lemma.  Therefore $L$ has two places above three;
each completion is $\mathbb Q_3(\zeta_3)$.
The field has no real places.  The Brauer invariant
sequence gives
\[
 \operatorname{Br}(\mathcal O_L[1/3])[3]
   =\{(a,-a):a\in\tfrac13\mathbb Z/\mathbb Z\}
   \simeq\mathbb F_3.                                      \tag{C.11.2}
\]
Kummer theory and the arithmetic etale comparison
surject $H^2(G_{L,S_L},\mu_3)$ onto this group.
These are the same comparison and invariant maps
used in A.14: Milne, II.2.9, II.2.1 and
Remark II.2.2(a) \cite{MilneADT}.  Since $\mu_3\subset L$,
fixing $1\mapsto\zeta_3$ identifies its coefficient with
the trivial module $\mathbb F_3$.  Thus
$H^2(G_{L,S_L},\mathbb F_3)\ne0$.

The image on $\mathcal K/3$ is a finite three-group,
because $L$ contains $\mathbb Q(\zeta_3)$ and the
action over that field has this property in A.1.
A nonzero module for a finite three-group over
$\mathbb F_3$ has a nonzero trivial quotient: its
coinvariants are nonzero by nilpotence of the
augmentation ideal.  Choose a surjection
$\mathcal K/3\twoheadrightarrow\mathbb F_3$.
The group $G_{L,S_L}$ has three-cohomological
dimension at most two.  The long exact sequence
therefore makes
\[
 H^2(G_{L,S_L},\mathcal K/3)
          \longrightarrow H^2(G_{L,S_L},\mathbb F_3)
\]
surjective, so its source is nonzero.

The degree-two finite-coefficient groups are finite.
Their inverse system is Mittag--Leffler, and the
finite-coefficient groups in degree three vanish.
The continuous-cohomology inverse-limit sequence
then gives $H^3_{\rm cts}(G_{L,S_L},\mathcal K)=0$.
Reduction modulo three in the long exact coefficient
sequence yields a surjection
\[
 H^2_{\rm cts}(G_{L,S_L},\mathcal K)/3
          \twoheadrightarrow H^2(G_{L,S_L},\mathcal K/3).
\]
This proves (C.11.1).

At every finite $v\nmid3$ the eigenvalues on
$\mathcal K(-1)$ have absolute values
$\#k(v)$ or $\sqrt{\#k(v)}$, by the boundary/elliptic
filtration.  None is one, so
$D_v=F_v-1$ is injective on this lattice.
The localization triangle of F.1 over $L$ thus has
zero degree-zero residue group.  Its long exact
sequence injects the nonzero finite-$S_L$ degree-two
group into the global degree-two group, with image
in the kernel of every displayed residue.
This proves the precise failure of the proposed
pure-residue replacement.
\end{proof}

\begin{statement}{C.12}{Proved}{The exact general number-field criterion with its retained arithmetic term}
For any number field $L$, put $S_L=\{v\mid3,\infty\}$,
$A_v=\mathcal K(-1)$ and $D_v=F_v-1$ for $v\notin S_L$.
Let $\mathcal R_L(A)$ denote the families in these
integral lattices tending to zero three-adically.
Then there are exact maps
\[
\begin{aligned}
 0\longrightarrow H^2_{\rm cts}(G_{L,S_L},\mathcal K)
 &\xrightarrow{\rm inf}H^2_{\rm cts}(G_L,\mathcal K)\\
 &\xrightarrow{\rm res}
       \mathcal R_L(A)/D\mathcal R_L(A)
       \longrightarrow0.
\end{aligned}                                              \tag{C.12.1}
\]
\[
 H^1_{\rm cts}(G_{L,S_L},\mathcal K)
 \xrightarrow{\sim}H^1_{\rm cts}(G_L,\mathcal K).             \tag{C.12.2}
\]
The quotient in (C.12.1) uses the actual image of $D$.
The first term is the inverse limit of the finite
groups $H^2(G_{L,S_L},\mathcal K/3^n)$.

For a section of $N^{\rm cc}$ over $L$, take its
actual local path-origin residues $r_v$ by the
formula (A.8.3), with residue size $\#k(v)$ and the
restricted $G_L$-action.  Its intrinsic obstruction
vanishes exactly when both of the following tests hold:
\[
 y_v=-D_v^{-1}r_v\in A_v\text{ for every }v\notin S_L,
        \qquad (y_v)_v\in\mathcal R_L(A);                  \tag{C.12.3}
\]
\[
 \eta(s)=0\quad\text{in }H^2_{\rm cts}(G_{L,S_L},\mathcal K).
                                                               \tag{C.12.4}
\]
Here, once (C.12.3) holds, $\eta(s)$ is the unique
class whose inflation is the full continuous
obstruction of $s$.  Every nonempty class-two lift
fiber is a torsor under
$H^1_{\rm cts}(G_{L,S_L},\mathcal K)$, with no further
geometric-stabilizer quotient.
This statement retains the finite-$S_L$ test; it
does not claim an evaluation of that test for
arbitrary $L$.
\end{statement}
\begin{proof}
Apply the finite-coefficient localization triangle
of F.1 and take its derived inverse limit.
The residue complex is the actual two-term complex
$[\mathcal R_L(A)\xrightarrow D\mathcal R_L(A)]$.
Each $D_v$ is injective by the weight calculation
at the end of C.11, so its degree-zero cohomology
is zero.  The finite-$S_L$ degree-one and degree-two
finite-coefficient groups are finite, and hence
their inverse systems are Mittag--Leffler.
Three-cohomological dimension at most two gives
$H^3_{\rm cts}(G_{L,S_L},\mathcal K)=0$ as in C.11.
The indicated pieces of the localization long
exact sequence are therefore precisely (C.12.1)
and its degree-one isomorphism; the same finite-group
Mittag--Leffler fact identifies the first term
with its finite-coefficient inverse limit.

The local section equation used in A.8 is unchanged
after restricting the arithmetic action to $G_L$:
the inertia comparison cochains are the restrictions
of the same based cusp paths, with the cusps now
permuted by $G_L$.  Applying that derivation over
each completion $L_v$ gives exactly $r_v$ as the
residue of the continuous obstruction class.
Its residue class vanishes precisely under
(C.12.3), because every $D_v$ is injective and
the sole integral correction is the displayed one.
Exactness of (C.12.1) then gives the unique
$\eta(s)$.  Its vanishing is equivalent to
vanishing of the original obstruction, proving
necessity and sufficiency of the two tests.

Once a lift exists, two lifts of a fixed
representative differ by a continuous
$\mathcal K$-valued one-cocycle; central
conjugacy changes it by a coboundary.
A geometric stabilizer of the underlying
$N^{\rm cc}$ section has fixed image in $V$,
which is zero by C.1.  Its remaining image is
in $\mathbb Z_3(1)$ and is again zero.
Thus no additional orbit identification occurs,
and the fiber is the asserted full torsor.
\end{proof}

\subsection{Algebraic evaluation of the point image}

\begin{statement}{C.16}{Proved}{Finite algebraic formulas for every point's abelian and central coordinates}
Put $m=3^n$.  Two successive Kummer covers and a quotient
of their finite deck group construct a based finite cover
$X_n\to U$ with group
\[
 P_n=N/(N^{3^n}\gamma_3N).                                  \tag{C.16.1}
\]
Its multiplication, arithmetic action, transition maps,
and fibers at algebraic points are computed by finite
algebraic operations.  The tangential base point is the
specified $s=z/x$, tangent one.

Choose compatible lifts $y_1,\ldots,y_{10}$ of a
$\mathbb Z_3$-basis of $V$.  For every $A\in U(L)$ the
construction gives compatible finite cocycles
\[
 c_{A,n}(g)=(v_{A,n}(g),z_{A,n}(g))
 \in(V/3^n)\oplus(\bigwedge\nolimits^2V/3^n).                \tag{C.16.2}
\]
They are obtained by evaluating roots of rational
functions at $A$ and comparing the resulting fibers
with the Puiseux fiber at $O$.  Their inverse limit
is a continuous cocycle whose class is the class-two
point section of $A$.  Thus (C.16.2) evaluates all
forty-five central coordinates with their tangent origin.
\end{statement}
\begin{proof}
We specify all the covers and comparison maps.
The geometric constructions are over
$\overline{\mathbb Q}$, and their finitely many
coefficients lie in computable number fields.

For a smooth affine curve $Y=C\setminus B$, with $C$
proper of genus $g$ and $B$ nonempty and finite, put
\[
 \mathcal F_m(Y)=
 \frac{\{f\in\overline{\mathbb Q}(C)^\times:
       \operatorname{ord}_x(f)\in m\mathbb Z
                         \text{ for every }x\in Y\}}
      {\overline{\mathbb Q}(C)^{\times m}}.                 \tag{C.16.3}
\]
Valuations give the exact sequence
\[
 0\longrightarrow\operatorname{Jac}(C)[m]
 \longrightarrow\mathcal F_m(Y)
 \longrightarrow
 \ker\bigl((\mathbb Z/m)^B\xrightarrow{\sum}\mathbb Z/m\bigr)
 \longrightarrow0.                                       \tag{C.16.4}
\]
The first map sends $[E]$ to the class of a function
with divisor $mE$.  An $m$th-power function makes
$E$ principal, proving injectivity.  If every boundary
valuation of $f$ is also divisible by $m$, its full
divisor is $mE$, proving exactness in the middle.
For surjectivity, lift a boundary vector to a degree-zero
divisor $D$ on $B$.  Choose $E$ with $m[E]=[D]$ in
the Jacobian and a function with divisor $D-mE$.
Multiplication by $m$ on the Jacobian is surjective
over $\overline{\mathbb Q}$, so these choices exist.
They give the required boundary valuations.

Choose a basis of $\operatorname{Jac}(C)[m]$ and the
boundary vectors $b-b_0$.  The preceding construction
gives $r_Y=2g+|B|-1$ functions $f_i$ forming a
$\mathbb Z/m$-basis of $\mathcal F_m(Y)$.  Normalize
$Y$ in
\[
 \overline{\mathbb Q}(C)
       (f_1^{1/m},\ldots,f_{r_Y}^{1/m}).                    \tag{C.16.5}
\]
Their independence gives degree $m^{r_Y}$.  At an
interior point remove the $m$-divisible valuation
from a Kummer equation; extracting a root of the
remaining unit is etale in characteristic zero.
Thus (C.16.5) is etale on $Y$.  Conversely every
cyclic etale cover of exponent dividing $m$ comes
from such a function, by function-field Kummer
theory and this valuation test.  This proves that
(C.16.5) is exactly the maximal geometric abelian
cover of exponent $m$.

All these are finite algebraic operations.  Represent
Jacobian classes by divisors, solve their finite
$m$-division equations, and compute the principal
functions by Riemann--Roch.  The latter computation
is given explicitly by Khuri-Makdisi's
Theorem/Algorithm 5.5 \cite{BC_KhuriMakdisi2004}.
Termination can also be certified without choosing
a Jacobian implementation: enumerate algebraic
rational functions, test the valuation condition
and their classes modulo $m$th powers, and retain
independent functions until their generated subgroup
has the known order $m^{r_Y}$ from (C.16.4).
Factoring $T^m-f$ over the algebraic function field
decides the power tests.  Equation (C.16.4) proves
that this enumeration reaches the required order.

Apply (C.16.5) to $U$, and denote the result by
$Y_m\to U$.  Its intrinsic maximal pro-three group
is $H_m=N^m[N,N]$.  Indeed the normal closure over
$U$ of a finite three-group cover of $Y_m$ again
has three-group Galois group: intersect its finitely
many conjugate subgroups, or use the associated
three-group wreath product.  This proves the
comparison with the open subgroup of $N$.
Apply (C.16.5) to $Y_m$, using its actual smooth
compactification and full boundary.  The composite
cover of $U$ has finite geometric group
\[
 E_m=N/(H_m^m[H_m,H_m]).                                   \tag{C.16.6}
\]
The denominator is characteristic in $N$, so this
composite is geometrically Galois.  Compute its
finite deck group and take the fixed field of
$E_m^m\gamma_3E_m$; this defines $X_n$.

In $N/(N^m\gamma_3N)$, the image of $H_m$ is
central and of exponent dividing $m$.  Hence
\[
 H_m^m[H_m,H_m]\subset N^m\gamma_3N.                        \tag{C.16.7}
\]
The preimage of $E_m^m\gamma_3E_m$ is therefore
exactly $N^m\gamma_3N$.  The quotient maps induced
by $N$ give inverse isomorphisms
$E_m/(E_m^m\gamma_3E_m)\simeq N/(N^m\gamma_3N)$.
This proves the precise map in (C.16.1).
Finite field automorphisms, fixed fields and the
maps of their normalizations are computed by
factoring finite algebraic equations.

Because $m$ is odd, the Hall map
\[
 \begin{split}
 (V/m)\oplus(\bigwedge\nolimits^2V/m)&\longrightarrow P_n,\\
 (v,z)&\longmapsto
 \prod_i y_i^{v_i}
 \prod_{i<j}[y_i,y_j]^{\,z_{ij}-v_iv_j/2}
 \end{split}                                              \tag{C.16.8}
\]
is an explicit isomorphism, with the products in
increasing order.  Its source has multiplication
\[
 (v,z)(v',z')=
       (v+v',z+z'+\tfrac12v\wedge v').                     \tag{C.16.9}
\]
Commuting successive generators verifies the map.
Conversely this source is generated by the ten
basis elements, has exponent $m$, and is of class
two.  The free-group map back is inverse on the
generators.  This proves both bijectivity and
$|P_n|=m^{55}$.  Abelian bases are read from the
Kummer pairing with (C.16.3), and their lifts are
found in the finite deck group.  They can be
chosen compatibly because $P_{n+1}\to P_n$ and
the maps on their abelian kernels are surjective.

Embed the cover fields compatibly in the Puiseux
field at $O$, using $s=z/x$ and fixed compatible
symbols $s^{1/3^a}$.  Galois acts on the algebraic
coefficients and fixes these symbols.
The characteristic subgroup in (C.16.1) is
preserved.  Equivalently there is a unique based
isomorphism $\phi_{g,n}:{}^gX_n\to X_n$ carrying
the conjugate chosen tangential fiber point to
the chosen one.  A deck transformation fixing a
fiber point is the identity.  Thus uniqueness
gives
$\phi_{gh,n}=\phi_{g,n}\,{}^g\phi_{h,n}$:
these are actual arithmetic descent maps.
They are computed from the finite function-field
isomorphisms and the Puiseux branch condition.
Compute coefficients until all finitely many
branches are distinguished; distinct algebraic
branches differ after finitely many terms.
The coefficient action agrees with the
path--Puiseux comparison of
\cite[\S1.1, (1.1.1)--(1.1.3)]{BC_HiranoMorishita2019},
applied here to the actual finite covers of $U$.

For $A\in U(L)$ choose compatible points $p_{A,n}$
in $X_n(A)$ by the root equations (C.16.5) and
the finite fixed-field maps.  The functions may
be chosen with valuation zero at every point
above $A$: multiply them by $m$th powers with
the required prescribed valuations, using
Riemann--Roch approximation.  This leaves their
classes in (C.16.3) unchanged.  Alternatively,
at a valuation $mk$, divide by a local parameter
to the power $mk$ and evaluate the resulting
unit; its root is trivialized by that parameter
to the power $k$.  These are explicit
isomorphisms of the Kummer equations.
The based inclusions associated to
$N^{3^{n+1}}\gamma_3N\subset N^{3^n}\gamma_3N$
supply the transition maps.  Their maps on
geometric fibers are surjective, allowing
successive compatible choices.

Let deck transformations act on the right on
fiber points, by composition with function-field
automorphisms.  Compute $c_{A,n}(g)$ from
\[
 \phi_{g,n}(g p_{A,n})=p_{A,n}\,c_{A,n}(g).                 \tag{C.16.10}
\]
All terms are algebraic points and maps already
constructed.  The descent identity gives
\[
 c_{A,n}(gh)=c_{A,n}(g)\rho_n(g)c_{A,n}(h).                 \tag{C.16.11}
\]
For fixed $n$, the equations, deck maps,
distinguishing Puiseux coefficients and fiber
points lie in a finite Galois extension of $L$.
Thus (C.16.10)--(C.16.11) are computed as
functions on a finite Galois group.
Formula (C.16.8) then extracts their coordinates.

Write the computed arithmetic action as
\[
 \rho_n(g)(v,z)=
       (R_gv,\bigwedge\nolimits^2R_g\,z+\ell_gv).
                                                               \tag{C.16.12}
\]
Here $\ell_g$ is read from the ten images
$\rho_n(g)y_i$ using (C.16.8).
It is linear because the class-two addition
$a+b=ab[a,b]^{-1/2}$ is recovered from the
group law; automorphisms preserve this
addition and the bracket.  In particular
the evaluated central cochain satisfies
\[
 \begin{split}
 z_{A,n}(gh)={}&z_{A,n}(g)
   +\bigwedge\nolimits^2R_g\,z_{A,n}(h)
   +\ell_g(v_{A,n}(h))\\
 &+\tfrac12v_{A,n}(g)\wedge R_gv_{A,n}(h).
 \end{split}                                              \tag{C.16.13}
\]
This retains the off-diagonal arithmetic term
as well as the tangential displacement.

Every finite quotient of $N/\gamma_3N$ has
finite three-power exponent, so it factors
through some $P_n$.  Therefore
$N/\gamma_3N\to\varprojlim_nP_n$ is an
isomorphism: compactness proves surjectivity,
and the kernel is tested in all finite
quotients.  The compatible finite cocycles
thus give a continuous cocycle.
The etale path comparison between the two
fiber functors on each $X_n$ is precisely
(C.16.10); cofinality identifies its inverse
limit with the point section.  A different
compatible fiber point changes the cocycle
by the corresponding compatible geometric
conjugation.  This proves the asserted class.
\end{proof}

\begin{statement}{C.17}{Proved}{An evaluated nonzero central point coordinate}
Let $\alpha^3=2$, $\zeta=\zeta_3$, $\eta=\zeta_9$
with $\eta^3=\zeta$, and $L=\mathbb Q(\zeta,\alpha)$.
For $A=[1:-\alpha:1]\in U(L)$ an explicit quotient
of $P_1$ of order $27$ has zero abelian point
coordinates and central coordinate
\[
                  -\kappa_L(\zeta)\in
                    H^1(G_L,\mathbb F_3)                  \tag{C.17.1}
\]
in the prescribed unit-tangent origin at $O$.
This class is nonzero.
\end{statement}
\begin{proof}
The functions $f=-x/z$ and $1-f=(x+z)/z$
are units on $U$: the line $x+z=0$ meets $E$
only at $Q_c$, with multiplicity three.
Pull back the cover
\[
 u^3=f,\qquad v^3=1-f,\qquad
 w^3=A(u):=(1-\zeta u)(1-\zeta^2u)^2.                      \tag{C.17.2}
\]
It is the $l=3,c=1$ Heisenberg cover of
\cite[\S2.1, (2.1.2), Theorem 2.1.4]
{BC_HiranoMorishita2019}.  We verify its
connected pullback and its deck maps.
The valuations of $f,1-f$ at $P,O$ are
$(1,0)$ and $(-1,-1)$, of determinant $-1$.
Their Kummer classes are therefore independent
and the first two equations have degree nine.
Since
\[
 \frac{A(\zeta u)}{A(u)}
                 =\frac{(1-u)^3}{1-u^3},                  \tag{C.17.3}
\]
the following are automorphisms:
\[
 \begin{array}{c|ccc}
       &u&v&w\\ \hline
 \tau  &\zeta u&v&(1-u)w/v\\
 \sigma&u&\zeta v&w\\
 z     &u&v&\zeta w .
 \end{array}                                              \tag{C.17.4}
\]
The identity $\prod_{j=0}^2(1-\zeta^ju)=1-u^3$
gives $\tau^3=1$; also $\sigma^3=z^3=1$
and $[\tau,\sigma]=z$.
If $A(u)=b^3$ in the degree-nine field, its
commuting automorphisms would satisfy
$\tau(b)/b=\zeta^i(1-u)/v$ and
$\sigma(b)/b=\zeta^j$ for some $i,j$.
Their two orders of application would then
differ by $\zeta$, a contradiction.
Thus the full degree is $27$.
Every zero and pole of these equations lies
over a deleted cusp, so the normalization
is etale on $U$.  Its exponent-three
class-two deck group is an actual quotient
of $P_1$.

At $A$ choose
\[
                 u=-1,\qquad v=\alpha,\qquad w=-\eta;
                                                               \tag{C.17.5}
\]
indeed $A(-1)=-\zeta$.
At $O$, put $r=s^{1/3}$.  The branches are
\[
 u=-r^{-1},\qquad
 v=r^{-1}(1+s)^{1/3},\qquad
 w=\eta^2r^{-1}(1+O(r)).                                   \tag{C.17.6}
\]
The last coefficient follows from
$r^3A(-r^{-1})|_{r=0}=\zeta^2$; the
remaining unit series has coefficients
in $\mathbb Q(\zeta)$.
For $g\in G_L$ with $g\eta=\zeta^e\eta$,
raw conjugation moves (C.17.5) by $z^e$
and the tangent branch by $z^{2e}$.
The based descent removes the latter
displacement.  Formula (C.16.10) therefore
gives $z^{e-2e}=z^{-e}$.
The raw arithmetic action commutes with
(C.17.4), and its based correction is
central, so the arithmetic action on the
finite group over $L$ is trivial.
Geometric conjugation fixes this central
cocycle.

Finally $X^3-2$ is Eisenstein at the
prime over two in $\mathbb Q(\zeta)$.
The distinct cyclic degree-three extension
$\mathbb Q(\eta)/\mathbb Q(\zeta)$ is
unramified there.  Consequently $\eta\notin L$
and $\kappa_L(\zeta)\ne0$.
The exact polynomial identities, the
twenty-seven deck elements and their
multiplication, both fiber displacements,
and the discriminants $-108$ and $-3^9$
are checked in
\texttt{code/bc\_point\_coordinates.py}
and its supplied output.
\end{proof}

\begin{statement}{C.18}{Proved}{The point-image set and its exact central membership criterion}
The point-image subset of the class-two germ set is
\[
 \left\{
   \left[\left(\varprojlim_n v_{A,n},
                   \varprojlim_n z_{A,n}\right)\right]:
                  A\in U(\overline{\mathbb Q})
 \right\},                                                \tag{C.18.1}
\]
where each finite function is algebraically evaluated
by (C.16.3)--(C.16.13).
This parametrization is injective, disjoint from the
nine packets (C.10.2), and equivariant for
$A\mapsto gA$.  Its stabilizer at $A$ is
$G_{\mathbb Q(A)}$.

Given a section and an algebraic point $A$, choose
a common number field $L$ of definition.
Let $d_n$ be the section cocycle in $P_n$ and
$c_{A,n}$ the evaluated cocycle above.
Their germs agree exactly when every finite set
\[
 \mathcal T_n=\{h\in P_n:
 d_n(g)=h^{-1}c_{A,n}(g)\rho_n(g)h
                         \text{ for all }g\in G_L\}
                                                               \tag{C.18.2}
\]
is nonempty.  Each test is a computation on a finite
Galois group.  Combined with the injective abelian
point map C.9, it gives a complete central
membership criterion for (C.18.1).
\end{statement}
\begin{proof}
C.16 proves (C.18.1), and C.9 proves injectivity,
disjointness and the stated stabilizer.
Conjugating the equations, divisors, roots and
based comparison maps of C.16 carries the
output at $A$ to the output at $gA$, exactly
as in (C.3), proving equivariance.

The equations in (C.18.2) factor through a
finite Galois extension containing the fields
of both finite cocycles and their actions.
They are therefore finite tests.
The inverse images of $\mathcal T_n$ in the
compact group $N/\gamma_3N$ are closed,
nonempty and nested.  Their intersection
provides one conjugator valid at every
precision, hence in $N/\gamma_3N$.
Conversely such a conjugator supplies
every finite solution.
Equality of germs is equality already over
$L$ by C.2.  This proves (C.18.2) without
assuming compatible choices of finite
conjugators.

We spell out the transport to the other
central coordinates.  The Kummer pairing
in (C.16.3) identifies $P_n^{\rm ab}$
with the actual $V/3^n$.
For cusp vectors it is the valuation
pairing; for proper vectors it is the
torsion-divisor pairing in (C.16.4).
Use the eight non-$O$ cusp vectors and
two lifts of an oriented Tate basis.
Actual cusp generators are obtained by
their Puiseux monodromy on the finite
covers.  Thus the basis comparison with
(E36.4) takes place in the actual deck
groups.  For two compatible generator
choices with the same abelian basis,
write $y_i'=y_i\exp(t_i)$, with
$t_i\in\bigwedge^2V$.  Substitution into
(C.16.8) gives the exact map from coordinates
in the primed generators to coordinates
in the unprimed generators:
\[
              (v,z)\longmapsto
                    (v,z+\sum_i v_i t_i).                 \tag{C.18.3}
\]
An abelian basis change $B$ gives $B$ on
$v$ and $\bigwedge^2B$ on central brackets,
followed by the computed linear correction
(C.18.3).

The based filling map $q_O$ gives
\[
 q_O(v,z)=(\pi v,\omega z+\lambda(v)),\qquad
               \omega=\bigwedge\nolimits^2\pi,            \tag{C.18.4}
\]
where $\lambda(v_i)$ is read from the
actual image $q_O(y_i)$ in the chosen
basis of $C$.  These images are computed
by the finite maps of covers.
One may take the two generators of $C$
to be $q_O(y_1),q_O(y_2)$ and the other
eight $y_i$ to be peripheral; then
$\lambda=0$.  The change to the
involution-compatible basis of C.13 is
(C.18.3), obtained by applying that
actual involution to the generators.
The orientation is
$k=[a,b]=-\delta$, as in (E36.5),
B.3 and A.8.

For the E36 parameter over a fixed
abelian cocycle $\alpha$, conjugate
the computed cocycle by an actual
geometric lift $\widetilde h$ of
$h\in V$, using
$c(g)\mapsto\widetilde h^{-1}c(g)
                       \rho(g)\widetilde h$.
Its abelian change is
$v(g)\mapsto v(g)+(R_g-1)h$.
The vector giving the prescribed
$\alpha$ is unique because
$H^0(G_L,V)=0$.
Write the prescribed $C$-reference as
$(\beta(g),w_\beta(g))$.
After this explicit comparison the
E36 parameter is the class of
\[
       \omega z_A(g)+\lambda(v_A(g))-w_\beta(g)
             \quad\text{in the generator }k.              \tag{C.18.5}
\]
The two $C$-cocycles have identical
abelian coordinate $\beta$; their
group difference is precisely this
central difference by (C.16.9).
Over a prescribed class-two lift of
the resulting $N^{\rm cc}$-cocycle,
the remaining fiber parameter is
likewise the difference of the
central cochains in
$\mathcal K=\ker\omega$, after
actual conjugation makes the
$N^{\rm cc}$-cocycles equal.
This proves that the construction
evaluates both the E36 parameter
and the forty-four-dimensional
fiber coordinate with their
specified reference translations.

C.2--C.3 identify the germs fixed
by $G_L$ with $\operatorname{Sec}_L$.
In particular
$\operatorname{Sec}_{\mathbb Q}$
is the entire $G_{\mathbb Q}$-fixed
subset.  Equations (C.10.2) and
(C.18.1) supply both geometric
image sets required in Task C.3.
The criterion ranges over all $n$;
it does not assert a bounded
finite stopping rule for equality
at all precisions.
\end{proof}

\begin{statement}{C.6}{Proved}{Complete class-two germ classification}
Task C.1 requests the full explicit residue complex,
obstruction residues, convergence criterion, and fibers
over every number field $L$.  C.11 disproves a
verbatim pure-residue extension of Part A; C.12
gives the exact replacement, including the
finite-$S_L$ arithmetic class and the lift torsor.
C.13 evaluates the preceding cuspidally central
obstruction by the local quadratic pairings at the
places above three, with their global sum relation.
C.14 computes the finite-$S_L$ complexes by finite Galois
cochain matrices and Artin--Verdier duality, and C.15
evaluates the retained class by explicit compact-support
cup traces.  Together with the restricted-family condition
in C.12, these give the full number-field criterion and
fibers requested in Task C.1.
For Task C.2, C.1--C.3 give injective restriction,
the actual Galois action and exact descent.  Formula
(C.5) identifies $\operatorname{Sec}_L$ with all
$G_L$-fixed germs.
For Task C.3, C.8--C.9 prove injectivity of the point
map and separation from cusp germs.  C.10 gives the
nine disjoint cusp parameter spaces with their torsion
and Galois action.  C.16 evaluates every point's
abelian and central coordinates by finite algebraic
operations at each precision; C.17 is an explicit
nonzero central instance.  C.18 gives the exact
point-image set, its equivariance and its full
central membership criterion.  Its coordinate
maps retain the E36 origin and the lift-fiber
translation.  Finally C.4--C.5 and C.7 prove
that both specified phantom classes remain
nonzero after every finite number-field
extension and in the germ colimit, settling
Task C.4.
\end{statement}
\begin{proof}
The cited statements provide the specified maps,
arithmetic computations and geometric image sets.
Their compatibility under finite extension is
the cochain restriction of C.2 and the based
fiber construction of C.16; hence they give
one classification over the filtered system
of number fields, as asserted.
\end{proof}

